\section{Purpose of the Project}
\label{sec:purpose}

\subsection{Goals}
\label{sec:goals}

The project generates the Kubernetes files of a small cluster from short
deployment documents. Each application's repository holds one \emph{project
document}, which states what the application needs: its processes, their
memory and CPU, the ports they offer, the services they depend on, and how the
application is reached from outside. One \emph{platform document} states what
all applications share, such as the entry points of the cluster, its probe
timings and its backup policies. Task~1 defined the two metamodels for this
work in Ecore. The \emph{source metamodel} describes both documents. The
\emph{target metamodel} describes the \emph{resolved model}, which records every
value the platform derives for an application, and the files that will be
generated from it.

Task~2 implements the model transformation between the two. It takes the
source models of every project in the cluster, the platform document, and three
documents that other processes publish, and produces the resolved model of one
project. The resolved model is the input of the code generation of Task~3. The
transformation is where the derivation rules of the specification become
executable: every value in a resolved model is computed by a rule in the
transformation, and none is written by an author.

The transformation has two goals.

\begin{enumerate}
  \item It derives the same decisions as the production implementation of the
    same specification, which is written in TypeScript and used to deploy the
    cluster. Both implementations are tested against the same committed result
    files, called \emph{oracles}.
  \item It stops with a message when it cannot derive a value. A resolved model
    with a missing or default value would generate files that look valid and
    deploy the wrong thing.
\end{enumerate}

\subsection{Role of the Transformation}
\label{sec:role}

Figure~\ref{fig:pipeline} shows the pipeline. Xtext parses each document into a
source model, and OCL rejects invalid ones (Task~1). A first QVT-Operational
transformation, the \emph{lowering}, copies values that an author declared once
for a whole project or application onto each process. Task~1 already contained
the lowering, because the constraints are checked on its output. The second
transformation, the \emph{resolution}, is the subject of this report. It reads
the lowered models and produces the resolved model. A small Java exporter then
writes the resolved dependency edges as JSON, which is compared with the
oracle. In Task~3, Acceleo templates will read the same resolved model and
write the Kubernetes files.

\begin{figure}[!htbp]
\centering
\begin{tikzpicture}[
  node distance=0.45cm and 0.55cm,
  box/.style={draw, rounded corners=2pt, align=center, font=\scriptsize,
    minimum height=0.9cm, text width=2.1cm},
  document/.style={box, fill=blue!6},
  stage/.style={box, fill=green!8},
  result/.style={box, fill=orange!10},
  arr/.style={-{Stealth[length=2mm]}}]
\node[document] (docs) {Project and platform documents (YAML)};
\node[document, below=of docs] (pinned) {Node contract, images lock, snapshot, migration proof};
\node[stage, right=of docs] (parse) {Xtext parsing and OCL constraints};
\node[stage, right=of parse] (lower) {Lowering (QVTo)};
\node[stage, right=of lower] (resolve) {\textbf{Resolution (QVTo)}};
\node[result, right=of resolve] (model) {Resolved model (XMI)};
\node[result, below=of model] (edges) {dependencies.json, compared with the oracle};
\node[stage, above=of model] (acceleo) {Code generation (Acceleo, Task~3)};
\draw[arr] (docs) -- (parse);
\draw[arr] (pinned) -| (parse);
\draw[arr] (parse) -- (lower);
\draw[arr] (lower) -- (resolve);
\draw[arr] (resolve) -- (model);
\draw[arr] (model) -- (edges);
\draw[arr] (model) -- (acceleo);
\end{tikzpicture}
\caption{The Pipeline, With the Resolution This Report Describes}
\label{fig:pipeline}
\end{figure}

\subsection{Refinement of Tasks 0 and 1}
\label{sec:refinement}

Task~1 planned one transformation from a source model, resolved against the
platform document, to a resolved model. Writing it showed that the platform
document is not enough. Three facts that the resolution needs come from other
documents:

\begin{itemize}
  \item the \emph{node contract}, which lists the machines of the cluster and
    what each can hold;
  \item the \emph{images lock}, which maps each image name in a project
    document to an immutable digest and the user the image runs as;
  \item the \emph{cluster-state snapshot}, which records on which machine each
    stored volume already lives.
\end{itemize}

A project may also publish a \emph{migration proof}, which records the
database schema version its new release was tested against, and it may name
configuration files, called \emph{Assets}, whose text ends up in the cluster.
The transformation therefore reads a third metamodel, the \emph{pinned
inputs}, beside the source metamodel. Appendix~\ref{app:metamodels} describes
this metamodel and the other changes to the Task~1 metamodels.

\subsection{Assumptions and Conditions}
\label{sec:assumptions}

\begin{itemize}
  \item The input is valid. The OCL constraints of Task~1 reject an invalid
    source model before the transformation runs. A value that only a set of
    documents can check, such as an image missing from the lock, stops the
    transformation instead.
  \item The transformation resolves one project at a time, but reads all
    projects of the cluster, because a dependency on another project's service
    needs that service's address and port.
  \item The production implementation is the reference. Where the
    specification leaves a choice, the transformation makes the same choice,
    so that both produce the same oracles.
  \item The transformation runs without Eclipse, through the standalone QVTo
    executor, so that the build server runs it on every change. A launch
    configuration is also provided to run it in Eclipse
    (Appendix~\ref{app:handin}).
\end{itemize}

\subsection{Scope}
\label{sec:scope}

The transformation covers every example project that the production
implementation resolves: \texttt{minimal}, \texttt{auth}, \texttt{data},
\texttt{delivery}, \texttt{edge}, \texttt{observability} and \texttt{secrets}.
Together these use every kind of value the specification derives: dependency
edges, secrets, volumes with backups, configuration files, sidecar containers,
database migrations, release gates and environment variables with
placeholders.

Two things are outside the scope of Task~2. The eighth example,
\texttt{knowledge}, is not yet resolved by the production implementation
either, so it has no reliable oracle. The target metamodel also does not yet
describe every file Task~3 must generate. The transformation produces the
generated-file entries that the target metamodel can describe, and the rest are
added with the templates in Task~3. Section~\ref{sec:limitations} discusses
both.
